\textbf{Commun à tous les candidats}

\medskip



























\textbf{Question 1} : 
On voit que pour $t = 5$, les coordonnées du point de la droite $\mathcal{D}'$ sont $(11~;~- 9~;~- 22)$ soit les coordonnées de M$_2$.

\smallskip

\hfill\textbf{Réponse b.}

\medskip

\textbf{Question 2} : Un vecteur directeur de la droite $\mathcal{D}'$ est : $\vect{u}_3\begin{pmatrix}3\\- 3\\-6\end{pmatrix}$.

\smallskip

\hfill\textbf{Réponse c.}

\medskip



\textbf{Question 3} : 






Un vecteur directeur de la droite $\mathcal{D}$ est $\vect{\text{AB}}\begin{pmatrix}-2\\2\\4\end{pmatrix}$ colinéaire au vecteur $\dfrac{1}{2}\vect{\text{AB}}\begin{pmatrix}-1\\1\\2\end{pmatrix}$

Un vecteur directeur de la droite $\mathcal{D}'$ est $\vect{u_3}\begin{pmatrix}3\\- 3\\-6\end{pmatrix}$ colinéaire au vecteur $\dfrac{1}{3}\vect{u_3}\begin{pmatrix}1\\- 1\\-2\end{pmatrix}$ ou encore colinéaire au vecteur $-\dfrac{1}{3}\vect{u_3}\begin{pmatrix}-1\\+ 1\\2\end{pmatrix}$.

Les droites $\mathcal{D}$ et $\mathcal{D}'$ ayant des vecteurs directeurs colinéaires au même vecteur sont donc parallèles.

De plus en remplaçant $t$ par $\dfrac{5}{3}$ dans l'équation paramétrique de $\mathcal{D}'$ on obtient $x = 1,~\, y = 1$ et $z = - 2$.

Les droites sont parallèles et ont un point commun : elles sont donc confondues.

\smallskip

\hfill\textbf{Réponse d.}

\medskip

\textbf{Question 4} : 





$\mathcal{P}$ a pour vecteur normal $\vect{p}\begin{pmatrix}1\\m\\-2\end{pmatrix}$.

$\mathcal{D}$ est parallèle au plan $\mathcal{P}$ si $\vect{\text{AB}}$ et $\vect{p}$ sont orthogonaux, soit :

$\vect{\text{AB}} \cdot \vect{p} = 0 \iff -2 \times 1 + 2 \times m + 4 \times (- 2) = 0 \iff - 2 + 2m - 8 = 0 \iff 2m = 10 \iff m = 5$. 

\smallskip

\hfill\textbf{Réponse c.}

\medskip

\bigskip
